\chapter{Allergen Immunotherapy (Allergy Shots)}

\section*{Overview}
Allergen immunotherapy, commonly known as allergy shots, is a treatment that gradually desensitizes a patient to specific allergens through repeated injections of gradually increasing doses. It is used to modify the immune response in allergic rhinitis, asthma, and stinging-insect hypersensitivity.

\section*{Alternative Names}
Allergy shots; allergen desensitization; subcutaneous immunotherapy (SCIT); hyposensitization therapy

\section*{Principle}
Repeated exposure to escalating allergen doses shifts the immune response from a type 1 hypersensitivity reaction driven by allergen-specific IgE toward tolerance, characterized by IgG-blocking antibodies, reduced mast cell and basophil reactivity, and induction of regulatory T cells. Over time, symptoms to natural allergen exposure are diminished. The treatment is administered in a buildup phase followed by a maintenance phase.

\section*{History}
Allergen immunotherapy was pioneered by Leonard Noon and John Freeman in 1911, who demonstrated that repeated grass pollen injections reduced allergic symptoms, and it has remained a cornerstone of allergy treatment for over a century.

\section*{Why This Test or Procedure Is Ordered}
Ordered for patients with allergic rhinitis or allergic asthma who have persistent symptoms despite optimal pharmacologic therapy and environmental control, and for those who cannot tolerate or wish to avoid long-term medication. It is also the treatment of choice for severe stinging-insect venom allergy and may be considered for atopic dermatitis triggered by aeroallergens.

\section*{Is This a Primary or Secondary Test or Procedure}
A secondary, disease-modifying therapeutic procedure used when allergen avoidance and pharmacotherapy are insufficient.

\section*{How This Test or Procedure Is Performed}
After skin or blood testing identifies the relevant allergens, the patient receives subcutaneous injections of a standardized allergen extract. Doses are increased weekly during a buildup phase of several months and then given at longer intervals during maintenance, typically every 2 to 4 weeks. Injections are always administered in a clinic where anaphylaxis can be treated.

\section*{What Preparation Is Needed}
The patient should have no acute illness or unstable asthma on the day of injection. Antihistamines may need to be held before skin testing, and the patient should remain in the clinic for at least 30 minutes after each injection for observation.

\section*{Contraindications and Precautions}
Unstable or poorly controlled asthma, significant cardiac disease, and inability to tolerate severe systemic reactions are relative contraindications. Immunotherapy should not be started during pregnancy, although maintenance may be continued, and beta-blockers are generally avoided because they impair treatment of anaphylaxis.

\section*{Risks}
Local swelling and redness at the injection site are common. Systemic reactions range from rhinitis, urticaria, and wheezing to rare, potentially life-threatening anaphylaxis, which is why injections require clinic observation.

\section*{What Happens After the Test or Procedure}
The patient is observed for at least 30 minutes after each injection and is advised on recognizing delayed reactions. Injections continue on a scheduled basis, with symptom diaries and medication use reviewed at each visit to guide dosing.

\section*{Understanding Your Results}
Improvement is measured by reduced symptom scores, decreased medication use, and improved quality of life, typically beginning within months of maintenance therapy. Venom immunotherapy may be confirmed by demonstrating tolerance through repeat sting challenge or by measuring venom-specific antibody levels.

\section*{Normal Results}
A significant reduction in allergic symptoms and medication requirement without systemic reactions during treatment.

\section*{Follow-up Tests and Procedures}
Symptoms, peak expiratory flow, and medication use are monitored at each visit, and skin or serum testing may be repeated to reassess sensitivity. After a successful course of 3 to 5 years, therapy is often discontinued and the patient followed for relapse.

\section*{References}
\begin{enumerate}
  \item MedlinePlus. Allergy shots. U.S. National Library of Medicine; 2025.
  \item Current Medical Diagnosis and Treatment. McGraw-Hill; 2025.
\end{enumerate}

\clearpage
